\documentclass[11pt]{article}
\usepackage[T1]{fontenc}
\usepackage[margin=1in]{geometry}
\usepackage{amsmath,amssymb,amsthm,booktabs,array}
\usepackage[colorlinks=true,urlcolor=blue,citecolor=blue,linkcolor=blue]{hyperref}
\newtheorem{theorem}{Theorem}
\newtheorem{proposition}[theorem]{Proposition}
\newcommand{\E}{\mathbb E}
\newcommand{\V}{V^{(k)}_{ij}}
\newcommand{\F}{\mathbb F_2}
\title{Butterfly gradient variances at half filling and four particles}
\author{Haobo Ma\thanks{ChronoAI Pte Ltd; The Omega Institute.}\and
Wenlin Zhang\thanks{National University of Singapore; The Omega Institute.}}
\date{18 October 2026}
\begin{document}
\maketitle
\begin{abstract}
For the single increasing-stride butterfly with uniform phase and rotation angles,
we calculate the first-phase gradient variance of the Walsh-spread magic-block input.
At half filling, $k=n/2$ and $n=4\cdot2^r$, an exact five-state transfer gives
$V^{(k)}_{ij}(n)/H_k(n)\sim(273/224)(9/8)^r$ for each of the six readouts
with $0\leq i<j\leq3$, where $H_k$ is the exact Haar variance.
Thus a comparison uniform over all even $2\leq k\leq n/2$ cannot have a finite
upper constant. A positive lower comparison does hold on this growing family
for $n\geq8$. We also complete the six $k=4$ readouts, proving their exact formulas,
the bounds $H_4/4<V^{(4)}_{ij}<H_4$, and their leading relative deficit.
The half-filling obstruction leaves the sparse regime $n-2k=\Omega(n)$ open.
\end{abstract}

\section{The scalar and the range of a uniform comparison}
We use the input and butterfly architecture of Kerenidis \cite{K}, keeping the
first phase on mode zero. The exact Haar scalar, rather than a relative-design
estimate, is the relevant benchmark for Claim 1, equation (40), of that paper.
In correspondence of 18 October \cite{communication}, Kerenidis suggested a
scalar comparison with constants uniform in particle number, reported finite-size
evidence up to half filling, and identified the recurrence family of two remaining
four-particle readouts. We verify the latter observation and show that half filling
requires a separate upper-bound statement. This does not refute the sparse
trainability assertion in Theorem 19 of \cite{K}, whose conclusion concerns
$k\leq\gamma_0 n$ for a sufficiently small constant $\gamma_0$.

All labels start at zero. Let $n=2^K$ and let $k$ be even, $2\leq k\leq n/2$.
Before Walsh spreading the state is
\[
 |\Psi_k\rangle=\bigotimes_{a=0}^{k/2-1}
 \frac{|1100\rangle+|0011\rangle}{\sqrt2}
 \otimes|0\rangle^{\otimes(n-2k)}.
\]
The normalized Walsh matrix is $H_n(a,c)=(-1)^{\operatorname{popcount}(a\mathbin{\&}c)}/\sqrt n$.
Write $\rho_{n,k}$ for the spread two-particle RDM on $\Lambda^2\mathbb C^n$,
normalized by $\operatorname{tr}\rho_{n,k}=\binom k2$.
The butterfly is $W_n=R_{K-1}D_{K-1}\cdots R_0D_0$, with independent
uniform angles modulo $2\pi$. Here $D_\ell$ is diagonal with entries
$e^{i\phi_{\ell,a}}$, and $R_\ell$ acts on pairs differing in bit $\ell$ by
$\left(\begin{smallmatrix}\cos\theta&-\sin\theta\\\sin\theta&\cos\theta\end{smallmatrix}\right)$.
In particular, $\E\cos^4\theta=3/8$; the two-mode blocks are not Haar distributed.
Set
\[
 L_{ij}=\langle i\wedge j|\Lambda^2(W_n)\rho_{n,k}\Lambda^2(W_n)^*|i\wedge j\rangle,
 \qquad V^{(k)}_{ij}(n)=\E(\partial_{\phi_{0,0}}L_{ij})^2.
\]
The derivative has mean zero by periodicity. A fixed diagonal data phase can be
absorbed into the first uniform phase layer.

The exact Haar identity of \cite[Theorem 9 and Appendix A.2]{K} is
\[
 H_k(n)=\frac{4\|B\|_F^2+2\|\Gamma(B)\|_F^2}{n^2(n^2-1)},\quad
 B=[d\Lambda^2(|0\rangle\langle0|),\rho_{n,k}],\quad
 \Gamma(B)_{ac}=\sum_d B_{ad,cd}.
\]
Pair indices in the contraction are extended antisymmetrically. For this input,
direct contraction gives
\begin{align}
 \|B\|_F^2&=\frac{k((k+2)n-2k^2+3k-6)}{2n^2},&
 \|\Gamma(B)\|_F^2&=\frac{k(k-1)^2(n-2k)}{n^2},\label{norms}\\
 H_k(n)&=\frac{2k((k^2-k+3)(n-2k)+7k-6)}{n^4(n^2-1)}.\label{haar}
\end{align}
For clarity, the first norm can be obtained before spreading: put
$u=n^{-1/2}(1,\ldots,1)$, $G=d\Lambda^2(|u\rangle\langle u|)$, and let $R$ be
the unspread RDM. Since $G$ is a projector,
$\|[G,R]\|_F^2=2(\operatorname{tr}(GR^2)-\operatorname{tr}(GRGR))$.
The two traces are $k(k+2)/(4n)$ and $k(2k^2-3k+6)/(4n^2)$, respectively.
The second norm in \eqref{norms} follows from
$\Gamma(R)=(k-1)P/2$, where $P$ projects onto the $2k$ active modes.
The vanishing of this second norm at $k=n/2$ explains why this endpoint must be
checked independently of fixed sparse filling.

\section{An exact growing family at half filling}
Put $\mathcal A=\{\{0,1\},\{0,3\},\{1,2\},\{2,3\}\}$ and
$\mathcal C=\{\{0,2\},\{1,3\}\}$.
\begin{theorem}\label{half}
Let $n=4\cdot2^r$, $r\geq0$, and $k=n/2$. For the six readouts above,
\begin{equation}
 n^4V^{(n/2)}_{ij}(n)=
 \begin{cases}
 \displaystyle\frac{273}{64}(9/8)^r-\frac72+
       \frac{535}{64}(9/16)^r-\frac{57}{8}(1/2)^r,&\{i,j\}\in\mathcal A,\\[2mm]
 \displaystyle\frac{273}{64}(9/8)^r-\frac72+
       \frac{161}{64}(9/16)^r-\frac{105}{32}(1/2)^r,&\{i,j\}\in\mathcal C.
 \end{cases}\label{half-formula}
\end{equation}
Consequently, for either family,
\[
 \frac{V^{(n/2)}_{ij}(n)}{H_{n/2}(n)}
 =\frac{273}{224}(9/8)^r-1+O((9/16)^r)\longrightarrow\infty.
\]
For every $r\geq1$ the ratio exceeds $1/4$.
\end{theorem}

Here is a finite transfer derivation, including the matrices needed to check the
formula without fitting values of $r$. Write $c=4h+c_0$, $u=4g+u_0$ and
$v=4(h+g)+v_0$, where high labels lie in $\F^r$ and low labels in $\{0,1,2,3\}$;
addition of labels means bitwise exclusive-or. Define
\[
 d_{ab,st}=\chi_a(s)\chi_b(t)-\chi_a(t)\chi_b(s),\quad
 \chi_a(s)=(-1)^{\operatorname{popcount}(a\mathbin{\&}s)},
\]
\[
 C_{ab,uv}=\frac12(d_{ab,01}+d_{ab,23})(d_{uv,01}+d_{uv,23})
             -\frac14\sum_{0\leq s<t<4}d_{ab,st}d_{uv,st}.
\]
At half filling the unspread RDM is, with $P_a$ projecting onto the four modes of block $a$,
\[
 R=\tfrac14 I_{\Lambda^2}+
   \sum_{a=0}^{n/4-1}(|\psi_a\rangle\langle\psi_a|-\tfrac14\Lambda^2P_a),
 \quad \psi_a=(e_{4a}\wedge e_{4a+1}+e_{4a+2}\wedge e_{4a+3})/\sqrt2.
\]
Walsh character orthogonality therefore gives, for $u,v\ne0$,
\begin{equation}
 (\rho_{n,n/2})_{0c,uv}
  =\frac{1}{4n}C_{0c_0,u_0v_0}\,
       \mathbf1_{h+g+v_{\rm high}=0}.\label{rdm-half}
\end{equation}
The nonzero $C_{0c_0,u_0v_0}$ are
$C_{01,01}=4$, $C_{01,23}=8$, $C_{02,02}=-4$, $C_{03,03}=-4$,
together with negatives on reversal of the second pair. Thus all nonzero
terms have $c_0+u_0+v_0=0$ and $c_0\ne0$.
Also $(\rho_{n,n/2})_{0c,cv}=0$ when $v\ne0$.

Write $x_a=(W_n)_{ia}$, $y_a=(W_n)_{ja}$ and $z_{ab}=x_ay_b-x_by_a$.
Differentiating $D_0$ gives $\partial L=i(T-\overline T)$, with
$T=\sum_{c>0,\,0<u<v}(\rho_{n,n/2})_{0c,uv}z_{0c}\overline{z_{uv}}$.
The first-layer phases give distinct frequencies for distinct disjoint triples;
the shared-index terms vanish by \eqref{rdm-half}. Consequently
\begin{equation}
 V^{(n/2)}_{ij}(n)=\frac1{16n^2}
 \sum_{\substack{c,u,v\ne0\text{ distinct}\\c+u+v=0}}
 C_{0c_0,u_0v_0}^{\,2}\,\E|z_{0c}|^2|z_{uv}|^2.\label{half-sum}
\end{equation}
The sum uses ordered $u,v$. Its factor accounts for both orientations of a
Fourier frequency and for $z_{vu}=-z_{uv}$.

Expand the four wedge factors in the expectation into sixteen terms, indexed
by a mask $m\in\{0,\ldots,15\}$. Their column labels are the four linear forms
$0,h,g,h+g$; masks reverse the order in each of
$(0,h),(g,h+g),(0,h),(g,h+g)$, the first two being ket factors.
When high bits are summed from the most significant bit downwards, equality of
the high suffixes is encoded by the kernel of $(a,b)\mapsto ah+bg$.
There are exactly five states, in the order
\[
 \F^2,\quad\langle(1,0)\rangle,\quad\langle(0,1)\rangle,
 \quad\langle(1,1)\rangle,\quad\{0\}.
\]
For a new bit choice $f\in(\F^2)^*$, the state changes from $S$ to
$S\cap\ker f$. Integrating the phases requires equal ket and bra counts in
each new coset, separately in each output row. For each old coset and row,
if $a$ factors stay and $b$ flip, the angular factor is
\begin{equation}
 a(a,b)=\E(\cos^a\theta\sin^b\theta)=
 \begin{cases}(2p)!(2q)!/(4^{p+q}p!q!(p+q)!),&a=2p,\ b=2q,\\0,&\text{otherwise}.
 \end{cases}\label{angular}
\end{equation}
For these high layers the two output rows use distinct gates, since $i,j<4$
already differ in a low bit. Flip signs disappear in every nonzero angular
integral. This gives the transfer independently of $r$.

After combining masks related by interchange of the middle two line states,
\eqref{half-sum} becomes
\begin{equation}
 n^4V^{(n/2)}_{ij}(n)=e^\mathsf T
 (M_A^r\ell_A+M_B^r\ell_B+M_C^r\ell_C),\qquad e=(1,0,0,0,0)^\mathsf T,\label{small-transfer}
\end{equation}
where
\[
 M_A=\frac1{16}\begin{pmatrix}9&1&9&1&0\\0&8&0&0&8\\0&0&18&0&2\\0&0&0&8&8\\0&0&0&0&16\end{pmatrix},\quad
 M_B=\frac1{16}\begin{pmatrix}9&1&0&0&0\\0&8&0&0&0\\0&0&0&0&0\\0&0&0&0&0\\0&0&0&0&0\end{pmatrix},
\]
\[
 M_C=\frac1{16}\begin{pmatrix}9&1&0&1&0\\0&8&0&0&0\\0&0&0&0&0\\0&0&0&8&0\\0&0&0&0&0\end{pmatrix}.
\]
The terminal vectors, obtained by integrating the two low-bit layers, are
{\small
\begin{equation}
\begin{array}{c|cc}
 &\mathcal A&\mathcal C\\\hline
 \ell_A&(9/8,53/8,49/64,49/16,7/2)^\mathsf T
        &(3/16,77/32,49/64,361/64,7/2)^\mathsf T\\
 \ell_B&(3/4,15/4,0,0,0)^\mathsf T
        &(-3/8,19/16,0,0,0)^\mathsf T\\
 \ell_C&(1/8,5/8,0,1/16,0)^\mathsf T
        &(3/16,13/32,0,41/64,0)^\mathsf T
\end{array}\label{terminals}
\end{equation}
}
For the five states, use the respective low-layer representatives
\[ (h,g)=(0,0),(0,1),(1,0),(1,1),(1,2). \]
Sum over $c_0,u_0$,
put $v_0=c_0+u_0$, weight by $C^2$, and use \eqref{angular} along the unique
paths through the first two layers. Exclude $u=0$ when $u_0=0$ and
$(0,1)\in S$, and exclude $v=0$ when $v_0=0$ and $(1,1)\in S$.
Mask signs are $(-1)^{\operatorname{popcount}(m)}$.
Masks $0,15$ give $M_A$ before combining with $5,10$; masks
$1,2,4,7,8,11,13,14$ give $M_B$; masks $3,12$ give $M_C$ before combining
with $6,9$. The matrices displayed are four times the moment transfer,
because $n^2/16=4^r$ in \eqref{half-sum}.
The third coordinate of $\ell_C$ is unreachable from $e$ and is set to zero.

\begin{proof}[Completion of the proof of Theorem \ref{half}]
The transfer derivation proves \eqref{small-transfer}. Direct multiplication gives
$P(M_A)=P(M_B)=P(M_C)=0$ for
\[
 P(x)=x(x-1)(x-1/2)(x-9/16)(x-9/8).
\]
Thus both sides of \eqref{half-formula} obey the same recurrence of order five.
Substitution of \eqref{terminals} in \eqref{small-transfer}, for $r=0,1,2,3,4$,
gives the displayed coefficients and proves the identities for every $r\geq0$.
This is an exact matrix calculation; it does not extrapolate a finite table.
Equation \eqref{haar} gives
\[
 n^4H_{n/2}(n)=\frac{n(7n/2-6)}{n^2-1}=\frac72+O(1/n),
\]
which proves the ratio asymptotic.
Let $F_C(r)$ be the second expression in \eqref{half-formula}. It is increasing:
\[
 F_C(r+1)-F_C(r)=\frac{273}{512}(9/8)^r
       -\frac{1127}{1024}(9/16)^r+\frac{105}{64}(1/2)^r>0.
\]
For $r\geq2$ the first term exceeds the magnitude of the second, and $r=0,1$
are direct rational checks. Since $F_C(1)=1099/1024$ and $n^4H_{n/2}<7/2$,
its ratio for $r\geq1$ exceeds $1099/3584>1/4$.
The first expression exceeds $F_C(r)$ by
$[187(9/16)^r-123(1/2)^r]/32>0$.
\end{proof}

\section{Completion of the four-particle readouts}
\begin{proposition}\label{four}
Put $n=8\cdot2^s$, $s\geq0$. Then
\[
 n^5V^{(4)}_{ij}(n)=\begin{cases}A_4(s),&\{i,j\}\in\mathcal A,\\
 C_4(s),&\{i,j\}\in\mathcal C,\end{cases}
\]
where
\begin{align*}
 A_4(s)&=120-\tfrac{891}{8}(3/4)^s+\tfrac{533}{64}(9/16)^s
       +\tfrac{179}{2}(1/2)^s-\tfrac{911}{8}(3/8)^s
       +\tfrac{8567}{128}(9/32)^s-40(1/4)^s,\\
 C_4(s)&=120-\tfrac{891}{8}(3/4)^s+\tfrac{533}{64}(9/16)^s
       -\tfrac{287}{4}(1/2)^s+\tfrac{125}{2}(3/8)^s
       +\tfrac{5281}{128}(9/32)^s-\tfrac{323}{8}(1/4)^s.
\end{align*}
For every one of these six readouts,
\[
 \frac14<\frac{V^{(4)}_{ij}(n)}{H_4(n)}<1,\quad
 H_4(n)=\frac{8(15n-98)}{n^4(n^2-1)},\quad
 1-\frac{V^{(4)}_{ij}(n)}{H_4(n)}
 =\frac{297}{320}(3/4)^s+O((9/16)^s).
\]
\end{proposition}

We give the exact recurrence behind the coefficient calculation. This also
specifies the computer-assisted part of this proposition. Set $q=8$, $m=2^s$,
and $R=\rho_{8,4}$. For rows $i,j$ put
$X_a=\sum_{h<m}(W_n)_{i,qh+a}$, $Y_a=\sum_{h<m}(W_n)_{j,qh+a}$,
$\alpha=(W_n)_{i0}$ and $\beta=(W_n)_{j0}$.
With $z_{ab}=X_aY_b-X_bY_a$ and $d_{ab}=\mathbf1_{a=0}(\alpha Y_b-X_b\beta)$
for $a<b$, differentiation gives
\[
 \partial L=im^{-2}T,\quad
 T=\sum_{a<b,\,u<v}R_{ab,uv}(d_{ab}\overline z_{uv}-z_{ab}\overline d_{uv}),
 \quad V=m^{-4}\E(-T^2).
\]
Independent subblocks $A,B$ at the next layer give
\[
 X'=uX^A+vX^B,\quad Y'=wY^A+zY^B,\quad
 \alpha'=u\alpha^A,\quad\beta'=w\beta^A,
\]
where $u=\cos\theta e^{i\eta_1}$, $v=-\sin\theta e^{i\eta_2}$,
$w=\cos\vartheta e^{i\eta_3}$ and $z=-\sin\vartheta e^{i\eta_4}$
are independent between rows, with uniform phases.
Order variables as $(X,Y,\overline X,\overline Y,\alpha,\beta,
\overline\alpha,\overline\beta)$ and write $M_s(e)=\E Z^e$.
Split the unmarked exponents between $f$ and $e-f$, putting every marked
exponent in $f$. Each split has coefficient
\[
 \prod_{\nu<4q}\binom{e_\nu}{f_\nu}\,
 a(d_X(f),d_X(e-f))a(d_Y(f),d_Y(e-f)),
\]
and contributes that coefficient times $M_s(f)M_s(e-f)$.
Here $d_X,d_Y$ count all factors in their row, including conjugates and marks.
Keep a split only if ket and bra degrees balance in each row of each branch,
and in each component column of each branch, counting the marks as column zero.
These rules follow from the independent new phases and from simultaneous shifts
of first-layer phases in each residue class. Initial moments are unique-path
phase and angle integrals in $W_8$.

Closing the surviving monomials of $-T^2$ under these splits gives 2,988 moments
and 10,003 transitions. The supplied rational certificate gives every exponent,
transition, base moment and exponential-polynomial solution for each readout.
The verifier checks each initial value and the entire coefficient identity
$M_{s+1}(e)=\sum_f a_{e,f}M_s(f)M_s(e-f)$, rather than a finite list of layers.
Its scalar contraction gives $A_4$ for $01,03,12,23$ and $C_4$ for $02,13$.
This proves the asserted all-layer formulas, including Kerenidis's two suggested
pairs $12,23$. The same certificate includes the preceding $q=4,k=2$ calculation:
for $n=4\cdot2^t$, all five pairs $01,02,12,13,23$ have relative deficit
$(17/20)(3/4)^t+O((9/16)^t)$. This is a fixed-particle effect, distinct from
Theorem \ref{half}.

For completeness, the bounds need no extrapolation. If $h_s=n^5H_4(n)$, then
$120-98\,2^{-s}\leq h_s<120$. We have $A_4(s)>C_4(s)$: their difference is
\[
 \tfrac{645}{4}2^{-s}-\tfrac{1411}{8}(3/8)^s
       +\tfrac{1643}{64}(9/32)^s+\tfrac38(1/4)^s>0.
\]
For $s\geq2$, dropping positive terms gives
$C_4(s)\geq120-\tfrac{891}{8}(3/4)^s-\tfrac{287}{4}2^{-s}
-\tfrac{323}{8}4^{-s}\geq2361/64>30$.
For $s\geq3$,
\[
 h_s-A_4(s)\geq(3/4)^s\left[
 \tfrac{891}{8}-\tfrac{375}{2}(2/3)^s
 -\tfrac{533}{64}(3/4)^s-\tfrac{8567}{128}(3/8)^s\right]>0;
\]
the bracket increases and is positive at $s=3$.
Direct rational substitution for $s=0,1,2$ completes the bounds.
The leading deficit follows by dividing the formulas by $h_s$.

\section{Exact finite-size comparison and the remaining question}
The tables in the verification archive cover every even $k\leq n/2$ at
$n=16,32,64,128$ for the representative retaining pair $01$ and disjoint pair $13$.
All entries are exact fractions. Table~\ref{numerics} reports the extrema,
rounded only for display; it makes no claim about other readout geometries.
\begin{table}[ht]
\centering
\begin{tabular}{rrrrr}
\toprule
$n$ & pair & minimum $V/H$ (at $k$) & maximum $V/H$ (at $k$)\\
\midrule
16 &01&0.729873 (4)&0.930905 (2)\\
16 &13&0.336871 (2)&0.597442 (8)\\
32 &01&0.588903 (8)&0.956252 (16)\\
32 &13&0.511095 (8)&0.787483 (16)\\
64 &01&0.536444 (16)&1.093098 (32)\\
64 &13&0.518345 (16)&0.991897 (32)\\
128 &01&0.516801 (32)&1.284232 (64)\\
128 &13&0.512561 (32)&1.223721 (64)\\
\bottomrule
\end{tabular}
\caption{Exact enumeration over the stated particle numbers, displayed as decimals.}
\label{numerics}
\end{table}

The computation uses a first-layer Fourier reduction independent of the
layer recurrence. For general $b=k/2$, define
$S_b(a)=\sum_{t<b}\chi_a(t)$. The spread matrix entries needed by the derivative are
\begin{align*}
 n^2(\rho_{n,k})_{0c,uv}
 &=C_{0c_0,u_0v_0}S_b((c+u+v)_{\rm high})\\
 &\quad+4\mathbf1_{u_0=0,c_0=v_0}S_b(u_{\rm high})S_b((c+v)_{\rm high})\\
 &\quad-4\mathbf1_{v_0=0,c_0=u_0}S_b(v_{\rm high})S_b((c+u)_{\rm high}).
\end{align*}
Disjoint triples have distinct first-layer frequencies; shared-index terms
are grouped by the frequency $\phi_0-\phi_v$. Integrating their products uses
the unique circuit paths and \eqref{angular}. This has a number of column sums
of order $n^3$, and the integrals do not depend on $k$. The archive includes
cross-checks against full spread RDMs, the all-layer certificate and the
independent five-state transfer.

The sparse uniform comparison is still open: neither these finite tables nor
the fixed-$k$ asymptotics prove constants for $k\leq\gamma n$ with fixed
$\gamma<1/2$. At half filling, the exact growing mode in Theorem \ref{half}
precludes a uniform upper constant. The neighborhood of half filling, where
the contraction term in \eqref{norms} becomes small, is a distinct remaining
boundary problem.

\paragraph{Computation and provenance.}
All certificate and integration checks use exact rational arithmetic; no Lean
formalization is claimed. The note and verification programs were developed
with Codex (GPT-6), with a supplied GPT Pro planning response. A separate
computational and mathematical referee pass checks the final claims.

\begin{thebibliography}{9}
\bibitem{K} I. Kerenidis, \emph{Scalable Quantum Machine Learning:
Trainability, Expressivity and Efficiency}, arXiv:2607.24014v2 [quant-ph], 2026.
\url{https://arxiv.org/abs/2607.24014v2}.
\bibitem{communication} I. Kerenidis, correspondence with H. Ma and W. Zhang,
18 October 2026 (readout-recurrence identification and uniform-scalar question).
\end{thebibliography}
\end{document}
